% GENERATED, edit the French unit or the translation table
\setcounter{chapter}{10}\mxunit{terminale-c09}{Orthogonality and distances}\label{ch:11}

{\small\sffamily\color{mxmint}S48, S51, S54, S56, S60\par}

\begin{keypoint}[Objectives]Calculate a dot product; justify orthogonality; find an orthogonal projection and a distance; calculate a volume.\end{keypoint}

\section{Course notes}
The coordinate calculations in this chapter are carried out in an orthonormal frame of space.
\subsection{Dot product in space}
For non-zero vectors $\vec u,\vec v$, $\vec u\cdot\vec v=\|\vec u\|\|\vec v\|\cos\theta$, where $\theta\in[0,\pi]$ is their angle. If one of them is zero, the dot product is $0$.
\[\vec u\cdot\vec v=\tfrac12(\|\vec u+\vec v\|^2-\|\vec u\|^2-\|\vec v\|^2).
\]
In an orthonormal frame, for $\vec u=(x,y,z)$ and $\vec v=(x',y',z')$,
\[\vec u\cdot\vec v=xx'+yy'+zz'.\]
The dot product is symmetric and bilinear; $\vec u\cdot\vec u=\|\vec u\|^2$. Two vectors are orthogonal if and only if their dot product is zero.
\paragraph{Example.} $\vec u=(1,2,2)$ and $\vec v=(2,-1,0)$ are orthogonal: $2-2+0=0$. Their norms are $3$ and $\sqrt5$.
\subsection{Line, plane and normal vector}
Two lines are orthogonal if their direction vectors are. A line is orthogonal to a plane if it is orthogonal to two non-collinear directions of that plane. A normal vector to a plane is a non-zero vector orthogonal to all its directions.
\paragraph{Why two directions suffice.} If $\vec n\cdot\vec u=\vec n\cdot\vec v=0$, then, for every $\vec w=s\vec u+t\vec v$, $\vec n\cdot\vec w=s\vec n\cdot\vec u+t\vec n\cdot\vec v=0$.
\begin{keypoint}[Vector condition for belonging to a plane]
If $A$ belongs to the plane $\mathcal P$ with normal vector $\vec n$, then $M\in\mathcal P$ if and only if $\vect{AM}\cdot\vec n=0$.
\end{keypoint}
\subsection{Orthogonal projection and distance}
The orthogonal projection $H$ of $M$ onto a plane is the unique point of the plane such that $\vect{MH}$ is normal to the plane; if $M$ belongs to the plane, $H=M$. The distance from $M$ to the plane is $MH$. On a line with direction vector $\vec u$, the projection satisfies $H\in d$ and $\vect{MH}\cdot\vec u=0$.
\paragraph{Proof of the minimum.} For every point $P$ of the plane, $\vect{MH}\perp\vect{HP}$, so $MP^2=MH^2+HP^2\ge MH^2$. Equality holds only for $P=H$. The projection is the closest point of the plane.
\paragraph{Method on a line.} If $d=A+\R\vec u$, write $H=A+t\vec u$ and solve $(A-M+t\vec u)\cdot\vec u=0$. Hence $t=\frac{\vect{AM}\cdot\vec{u}}{\lVert\vec u\rVert^2}$.
\subsection{Angles, areas and volumes}
For an angle $\widehat{ABC}$, use $\vect{BA}$ and $\vect{BC}$: $\cos \widehat{ABC}=\frac{\vect{BA}\cdot\vect{BC}}{BA\times BC}$. For the acute angle between two lines, take the absolute value of the quotient.
The area of a triangle is $\frac{bh}{2}$; the volume of a tetrahedron is $\frac{\mathcal{A}_{\rm base}h}{3}$, where $h$ is an orthogonal distance. The units are squared and cubed respectively.
\subsection{Expansions and loci}
By bilinearity,
\[\|\vec u+\vec v\|^2=\|\vec u\|^2+2\vec u\cdot\vec v+\|\vec v\|^2,\qquad
\|\vec u-\vec v\|^2=\|\vec u\|^2-2\vec u\cdot\vec v+\|\vec v\|^2.\]
We can therefore also calculate $\vec u\cdot\vec v=\frac12(\|\vec u\|^2+\|\vec v\|^2-\|\vec u-\vec v\|^2)$.
For two distinct points $A,B$ with midpoint $I$, the locus of points $M$ such that $MA=MB$ is the plane through $I$ with normal vector $\vect{AB}$: the \emph{perpendicular bisector plane}. Indeed, writing the vectors from $I$, the expansions give $MA^2-MB^2=2\vect{IM}\cdot\vect{AB}$. Equality of the distances is therefore equivalent to $\vect{IM}\cdot\vect{AB}=0$.
\textbf{Extension, not assessed in 2026--2027.} Two planes are perpendicular when their normal vectors are orthogonal. This property must not be confused with parallelism of planes, which corresponds to collinear normal vectors.


\section{Research activity}\textbf{The point of a line closest to a given point --- 35 min.}\par In an orthonormal frame of space, consider $A(1;0;0)$, $M(0;2;3)$ and the line $d$ through $A$ with direction vector $\vec u=(1;1;0)$. We look for the point of $d$ closest to $M$.
\begin{enumerate}\item Express, in terms of a real parameter $t$, the coordinates of a point $H$ of $d$.
\item Find the value of $t$ for which $\vect{MH}\cdot\vec u=0$. Deduce the coordinates of the resulting point $H$.
\item Calculate the distance $MH$.
\item Let $P$ be any point of $d$. Prove that $MP^2=MH^2+HP^2$, then deduce that $H$ is the unique point of $d$ closest to $M$.\end{enumerate}

\section{Progressive exercises}

\mxexerciseitem{1}{Practice}In an orthonormal basis of space, consider $\vec u=(2;-1;3)$ and $\vec v=(1;4;0)$. Calculate $\vec u\cdot\vec v$.

\mxexerciseitem{2}{Practice}In an orthonormal basis, consider $\vec u=(a;2;1)$ and $\vec v=(3;-1;4)$, where $a$ is real. Find $a$ so that $\vec u$ and $\vec v$ are orthogonal.

\mxexerciseitem{3}{Practice}In an orthonormal frame, a plane is spanned by $\vec u=(1;-1;0)$ and $\vec v=(1;0;-1)$. Show that $\vec n=(1;1;1)$ is normal to it. Then consider $A(0;0;0)$ and $B(2;2;2)$. Describe geometrically the set of points equidistant from $A$ and $B$.

\mxexerciseitem{4}{Practice}In an orthonormal basis, consider $\vec u=(1;0;1)$ and $\vec v=(0;1;1)$. Find the measure, between $0$ and $\pi$, of the angle between these two vectors.

\mxexerciseitem{5}{Practice}A tetrahedron has as its base a right-angled triangle whose two perpendicular sides measure $3$ cm and $4$ cm. The height of the tetrahedron relative to this base is $5$ cm. Calculate its volume in $\mathrm{cm}^3$.

\section{Textbook or past-paper work}\textbf{Barbazo, p. 362, exercises 23 and 27.} In S48: exercise 23, to calculate a dot product in coordinates. At home after S54: exercise 27; check orthogonality to the two direction vectors of a plane, without settling for a single zero product.

\section{Homework}\textbf{HW11 --- A face and its height --- 50 min.}\par In an orthonormal frame with origin $O$, consider the points $A(2;0;0)$, $B(0;2;0)$ and $C(0;0;2)$.
\begin{enumerate}\item Prove that the triangle $ABC$ is equilateral and calculate its area.
\item Show that the vector $\vec n=(1;1;1)$ is normal to the plane $(ABC)$.
\item Let $H$ be the orthogonal projection of $O$ onto $(ABC)$. Justify that its coordinates are of the form $(t;t;t)$. Find $t$ using $\vect{AH}\cdot\vec n=0$.
\item Calculate $OH$, then the volume of the tetrahedron $OABC$ taking $ABC$ as the base. Check this volume by taking $OAB$ as the base.\end{enumerate}

\ifmxteacher\mxnewpage\section{Teaching notes and solutions}\begin{teachernote}[Teacher notes]Begin with the three expressions of the dot product. Orthogonal lines are not necessarily intersecting: ``perpendicular'' requires an intersection. To prove that a line is orthogonal to a plane, choose two non-collinear directions. The general point-to-plane distance formula is not required by heart: favour the orthogonal projection. The Cartesian equations used in the S54 progression can be introduced through the vector condition, then formalised in S67.\end{teachernote}

\subsection{Solutions to the activity}\begin{enumerate}
\item A point $H$ of $d$ is written $H=A+t\vec u$, so $H(1+t,t,0)$ for $t\in\R$.
\item $\vect{MH}=(1+t,t-2,-3)$. The dot product with $\vec u=(1,1,0)$ is $1+t+t-2=2t-1$. Setting it to zero gives $t=\frac{1}{2}$ and $H(\frac{3}{2},\frac{1}{2},0)$.
\item The frame is orthonormal: $MH^2=(\frac{3}{2})^2+(-\frac{3}{2})^2+(-3)^2=\frac{27}{2}$. So $MH=\frac{3\sqrt{6}}{2}$.
\item For $P\in d$, $\vect{HP}$ is collinear with $\vec u$ and hence orthogonal to $\vect{MH}$. Pythagoras' theorem gives $MP^2=MH^2+HP^2\ge MH^2$, with equality only if $P=H$. The point found is indeed the unique point of the line closest to $M$.
\end{enumerate}

\paragraph{E1.}In the orthonormal frame, the dot product is $2\times1+(-1)\times4+3\times0=-2$.

\paragraph{E2.}The dot product is $3a+2(-1)+1\times4=3a+2$. Orthogonality is equivalent to its vanishing: $3a+2=0$, so $a=-\frac{2}{3}$.

\paragraph{E3.}The two direction vectors are non-collinear and both dot products are zero. The midpoint is $I(1,1,1)$. The required locus is the plane through $I$ with normal vector $\vect{AB}=(2,2,2)$, hence with normal vector $(1,1,1)$.

\paragraph{E4.}Both vectors are non-zero, with norms $\sqrt2$, and their dot product is $1$. So their angle $\theta\in[0,\pi]$ satisfies $\cos \theta=\frac{1}{\sqrt2\sqrt2}=\frac{1}{2}$, that is $\theta=\frac{\pi}{3}$.

\paragraph{E5.}The area of the right-angled base is $\frac{3\times4}{2}=6\,\mathrm{cm}^2$. The height perpendicular to this base is $5$ cm, so $V=\frac{6\times5}{3}=10\,\mathrm{cm}^3$.

\subsection{Solutions to HW11}\begin{enumerate}
\item $\vect{AB}=(-2,2,0)$, $\vect{AC}=(-2,0,2)$ and $\vect{BC}=(0,-2,2)$; the three norms are $2\sqrt2$. The triangle is equilateral. Its height is $\sqrt{(2\sqrt2)^2-(\sqrt2)^2}=\sqrt6$; its area is therefore $\frac{(2\sqrt2)(\sqrt6)}{2}=2\sqrt{3}$.
\item The vectors $\vect{AB},\vect{AC}$ are non-collinear. Their dot products with $(1,1,1)$ are respectively $-2+2=0$ and $-2+2=0$. The non-zero vector $\vec n$ is therefore normal to $(ABC)$.
\item The projection of $O$ lies on the line through $O$ with direction vector $\vec n$, so $H=(t,t,t)$. Its condition for belonging to the plane is $(t-2)+t+t=0$, that is $t=\frac{2}{3}$. We obtain $H(\frac{2}{3},\frac{2}{3},\frac{2}{3})$, which satisfies both conditions for an orthogonal projection.
\item $OH=\sqrt{3(\frac{2}{3})^2}=\frac{2}{\sqrt{3}}$. Taking $ABC$ as the base, the volume is $V=\frac{(2\sqrt3)(\frac{2}{\sqrt{3}})}{3}=\frac{4}{3}$ volume units. Other base: the triangle $OAB$ is right-angled at $O$, with area $\frac{2\times2}{2}=2$; the line $(OC)$ is perpendicular to the plane $(OAB)$, and $OC=2$. Thus $V=\frac{2\times2}{3}=\frac{4}{3}$.
\end{enumerate}

\subsection{Solutions to the textbook or past-paper work}\paragraph{Exercise 23.} In the given orthonormal frame, the dot product is the sum of the products of the corresponding coordinates:
\[(-3.5)(-4)+2\times5.5+(-4)(-3)=14+11+12=37,\]
\[\frac23\left(-\frac34\right)+\frac{15}7\frac{14}3+\left(-\frac{12}{11}\right)\left(-\frac{22}4\right)=-\frac12+10+6=\frac{31}2.\]
\paragraph{Exercise 27.} In each case, the two given vectors $\vec u$ and $\vec v$ span the plane and are not collinear. A non-zero vector $\vec w$ is normal if, and only if, it is orthogonal to each of them.
\begin{enumerate}
\item $\vec w\cdot\vec u=-132+312-180=0$ and $\vec w\cdot\vec v=-66+546-480=0$. The vector $\vec w$ is normal to the plane.
\item $\vec w\cdot\vec u=20+6-26=0$, but $\vec w\cdot\vec v=18-6-10=2\ne0$. It is not normal to the plane.
\item $\vec w\cdot\vec u=\frac12+6-7=-\frac12\ne0$. This single non-zero product is enough to rule out normality.
\end{enumerate}

\fi

